\documentclass[aspectratio=169,11pt]{beamer}

\usetheme{Madrid}
\usecolortheme{seahorse}
\usepackage[utf8]{inputenc}
\usepackage[T1]{fontenc}
\usepackage[brazil]{babel}
\usepackage{booktabs}
\usepackage{tikz}
\usetikzlibrary{positioning,arrows.meta}
\usepackage{listings}

\lstset{
  basicstyle=\ttfamily\scriptsize,
  breaklines=true,
  columns=fullflexible,
  frame=single,
  showstringspaces=false,
  literate={á}{{\'a}}1 {ã}{{\~a}}1 {ç}{{\c{c}}}1 {é}{{\'e}}1 {í}{{\'i}}1 {ó}{{\'o}}1 {ú}{{\'u}}1 {Á}{{\'A}}1
}

\title{Comparação de tecnologias de invocação de serviços remotos}
\subtitle{Serviço de streaming de músicas --- beta, sem teste de carga}
\author{Integrante 1 \and Integrante 2 \and Integrante 3}
\institute{Computação Distribuída\\Prof. Nabor C. Mendonça}
\date{Novembro de 2026}

\begin{document}

\begin{frame}
  \titlepage
\end{frame}

\begin{frame}{Equipe}
  \begin{itemize}
    \item Integrante 1 --- preencher nome e matrícula
    \item Integrante 2 --- preencher nome e matrícula
    \item Integrante 3 --- preencher nome e matrícula
  \end{itemize}
  Linguagem única do trabalho: Python 3.\\
  Este arquivo é um beta: a comparação de carga fica para a versão final.
\end{frame}

\begin{frame}{Roteiro}
  \tableofcontents
\end{frame}

\section{Serviço}

\begin{frame}{O que o enunciado pede}
  Quatro tecnologias de invocação remota cobrem a mesma necessidade: SOAP, REST, GraphQL e gRPC.
  \vspace{0.4em}
  O serviço de streaming gerencia três recursos ligados entre si:
  \begin{itemize}
    \item usuários;
    \item músicas;
    \item playlists.
  \end{itemize}
  O cliente precisa listar todos os usuários, todas as músicas, as playlists de um usuário, as músicas de uma playlist e as playlists que contêm uma música.
\end{frame}

\begin{frame}{Modelo do diagrama}
  \centering
  \begin{tikzpicture}[
      caixa/.style={draw, align=center, font=\small, inner sep=3pt, minimum width=2.7cm},
      seta/.style={-{Stealth[length=2mm]}, font=\scriptsize}
    ]
    \node[caixa] (usuarios) {coleção\\Usuários};
    \node[caixa, below=0.85cm of usuarios] (usuario) {Usuário\\{\scriptsize ID, Nome, Idade}};
    \node[caixa, right=1.15cm of usuario] (playlists) {coleção\\Playlists};
    \node[caixa, below=0.85cm of playlists] (playlist) {Playlist\\{\scriptsize ID, Nome}};
    \node[caixa, right=1.15cm of playlists] (musicas) {coleção\\Músicas};
    \node[caixa, below=0.85cm of musicas] (musica) {Música\\{\scriptsize ID, Nome, Artista}};
    \draw[seta] (usuarios) -- node[right] {0..*} (usuario);
    \draw[seta] (musicas) -- node[right] {0..*} (musica);
    \draw[seta] (usuario) -- node[above] {1..1} (playlists);
    \draw[seta] (playlists) -- node[right] {0..*} (playlist);
    \draw[{Stealth}-{Stealth}, font=\scriptsize] (playlist.east)
      -- node[near start, above] {0..*} node[near end, above] {1..*} (musica.west);
  \end{tikzpicture}
\end{frame}

\begin{frame}{Regras que saem do diagrama}
  \begin{itemize}
    \item O catálogo guarda 0..* usuários e 0..* músicas.
    \item Cada usuário tem exatamente uma coleção de playlists, e essa coleção tem 0..* playlists.
    \item Uma playlist contém 1..* músicas. Uma música entra em 0..* playlists.
    \item Por isso a playlist já nasce com ao menos uma música e a API recusa esvaziá-la.
    \item Apagar o usuário apaga as playlists dele. As músicas continuam no catálogo.
  \end{itemize}
  A implementação guarda os dados em memória. Reiniciar o processo volta à carga inicial: Ana, Bruno, quatro músicas e três playlists.
\end{frame}

\begin{frame}{As cinco consultas}
  \small
  \begin{tabular}{@{}lll@{}}
    \toprule
    Consulta & REST & GraphQL / SOAP / gRPC \\
    \midrule
    Todos os usuários & \texttt{GET /usuarios} & \texttt{usuarios} / \texttt{listarUsuarios} \\
    Todas as músicas & \texttt{GET /musicas} & \texttt{musicas} / \texttt{listarMusicas} \\
    Playlists do usuário & \texttt{GET /usuarios/\{id\}/playlists} & \texttt{playlistsDoUsuario} \\
    Músicas da playlist & \texttt{GET /playlists/\{id\}/musicas} & \texttt{musicasDaPlaylist} \\
    Playlists com a música & \texttt{GET /musicas/\{id\}/playlists} & \texttt{playlistsComMusica} \\
    \bottomrule
  \end{tabular}

  \vspace{0.7em}
  SOAP e gRPC usam os mesmos nomes de operação: \texttt{listarPlaylistsDoUsuario}, \texttt{listarMusicasDaPlaylist} e \texttt{listarPlaylistsComMusica}.
\end{frame}

\section{Tecnologias}

\begin{frame}{Como o código está organizado}
  As quatro APIs chamam o mesmo núcleo, \texttt{streaming/servico.py}.
  A regra de negócio não muda com o protocolo. O que muda é o contrato, o formato da mensagem e o estilo da chamada.
  \vspace{0.5em}
  \begin{tabular}{@{}ll@{}}
    \toprule
    Servidor & Endereço \\
    \midrule
    REST & \texttt{127.0.0.1:8001} \\
    GraphQL & \texttt{127.0.0.1:8002/graphql} \\
    SOAP & \texttt{127.0.0.1:8003/?wsdl} \\
    gRPC & \texttt{127.0.0.1:8004} \\
    \bottomrule
  \end{tabular}
\end{frame}

\begin{frame}{SOAP}
  \textbf{Origem.} O SOAP 1.1 surgiu em 2000, com a Microsoft e parceiros. O W3C padronizou o SOAP 1.2. O nome deixou de ser lido como sigla.
  \vspace{0.35em}
  \textbf{Características.} Envelope XML, falhas explícitas e contrato WSDL. O transporte usual é HTTP. A chamada tem nome de operação, no estilo RPC.
  \vspace{0.35em}
  \textbf{Vantagens.} Contrato formal, ferramentas que geram cliente a partir do WSDL e modelo de erro padronizado.
  \vspace{0.35em}
  \textbf{Desvantagens.} XML verboso, difícil de cachear e trabalhoso para um CRUD simples. Cada consulta do enunciado virou uma operação a mais no WSDL.
\end{frame}

\begin{frame}[fragile]{SOAP --- trecho}
\begin{lstlisting}[language=Python]
class StreamingSoap(ServiceBase):
    @rpc(Unicode, Integer, _returns=UsuarioSoap)
    def criarUsuario(ctx, nome, idade):
        return _chamar(lambda: _usuario(
            servico.criar_usuario(nome, idade)))

    @rpc(Integer, _returns=Array(PlaylistSoap))
    def listarPlaylistsDoUsuario(ctx, usuarioId):
        return _chamar(lambda: [_playlist(item) for item
            in servico.listar_playlists_do_usuario(usuarioId)])
\end{lstlisting}
O WSDL sai em \texttt{http://127.0.0.1:8003/?wsdl}. Um erro de negócio volta como \texttt{Fault} com código \texttt{Client.NAO\_ENCONTRADO}, \texttt{Client.VALIDACAO} ou \texttt{Client.CONFLITO}.
\end{frame}

\begin{frame}{REST}
  \textbf{Origem.} REST é o estilo arquitetural descrito por Roy Fielding na tese de 2000. Não é um protocolo nem um formato.
  \vspace{0.35em}
  \textbf{Características.} Recursos identificados por URL, verbos HTTP e representação JSON neste trabalho. A aplicação é sem estado: cada pedido traz o que o servidor precisa.
  \vspace{0.35em}
  \textbf{Vantagens.} Mapeia direto o CRUD das três coleções. É simples de testar no navegador e em \texttt{curl}. O cache do HTTP cabe melhor aqui do que nas outras três.
  \vspace{0.35em}
  \textbf{Desvantagens.} As consultas de relacionamento pedem rotas extras. O cliente que quer playlist e músicas juntas depende de como o servidor montou o JSON. O formato de erro é convenção da API, não do estilo.
\end{frame}

\begin{frame}[fragile]{REST --- trecho}
\begin{lstlisting}[language=Python]
@app.post("/usuarios", status_code=201)
def criar_usuario(corpo: UsuarioIn):
    return servico.criar_usuario(corpo.nome, corpo.idade)

@app.get("/usuarios/{usuario_id}/playlists")
def playlists_do_usuario(usuario_id: int):
    return servico.listar_playlists_do_usuario(usuario_id)

@app.get("/playlists/{playlist_id}/musicas")
def musicas_da_playlist(playlist_id: int):
    return servico.listar_musicas_da_playlist(playlist_id)
\end{lstlisting}
A documentação interativa fica em \texttt{/docs}. \texttt{DELETE} responde 204. Validação responde 400, ausência responde 404 e música repetida na playlist responde 409.
\end{frame}

\begin{frame}{GraphQL}
  \textbf{Origem.} O Facebook criou o GraphQL em 2012 e o abriu em 2015. Hoje está na GraphQL Foundation.
  \vspace{0.35em}
  \textbf{Características.} Um endpoint, um schema e uma consulta em que o cliente escolhe os campos. Há consultas e mutações. O servidor deste trabalho publica o schema em \texttt{/graphql}, com GraphiQL.
  \vspace{0.35em}
  \textbf{Vantagens.} As cinco consultas do enunciado cabem no schema, e o cliente ainda pode navegar de usuário para playlists e de música para playlists no mesmo pedido. Dá para pedir só \texttt{nome}, sem a lista de músicas.
  \vspace{0.35em}
  \textbf{Desvantagens.} Cache HTTP fica mais difícil. O servidor resolve campos sob demanda, então uma consulta aninhada repete idas ao núcleo. A curva é maior que a do REST para um CRUD tão pequeno.
\end{frame}

\begin{frame}[fragile]{GraphQL --- trecho}
\begin{lstlisting}
query ConsultasDoEnunciado {
  usuarios { id nome idade }
  musicas { id nome artista }
  playlistsDoUsuario(usuarioId: 1) {
    nome
    musicas { nome artista }
  }
  musicasDaPlaylist(playlistId: 1) { nome artista }
  playlistsComMusica(musicaId: 2) { id nome }
}
\end{lstlisting}
A mutação \texttt{criarPlaylist} exige \texttt{nome}, \texttt{usuarioId} e \texttt{musicaIds}. Lista vazia é recusada pelo núcleo, com a mesma mensagem das outras APIs.
\end{frame}

\begin{frame}{gRPC}
  \textbf{Origem.} O gRPC foi aberto pelo Google em 2015, a partir do sistema interno Stubby. Usa HTTP/2 e Protocol Buffers.
  \vspace{0.35em}
  \textbf{Características.} O contrato é o arquivo \texttt{.proto}. As ferramentas geram o servidor e o cliente. A mensagem é binária. A chamada tem nome, como no SOAP, com outro formato.
  \vspace{0.35em}
  \textbf{Vantagens.} Contrato rígido, mensagens compactas e bom encaixe entre serviços. O mesmo arquivo descreve as vinte operações do CRUD e das consultas.
  \vspace{0.35em}
  \textbf{Desvantagens.} Exige gerar código antes de subir o servidor. A mensagem binária é ruim de inspecionar à mão. No protobuf, o valor ausente e o zero se confundem, então nome e idade na atualização são \texttt{optional}.
\end{frame}

\begin{frame}[fragile]{gRPC --- trecho}
\begin{lstlisting}
service StreamingMusicas {
  rpc CriarUsuario(CriarUsuarioRequest) returns (Usuario);
  rpc ListarPlaylistsDoUsuario(Id) returns (ListaPlaylists);
  rpc ListarMusicasDaPlaylist(Id) returns (ListaMusicas);
  rpc ListarPlaylistsComMusica(Id) returns (ListaPlaylists);
}

message AtualizarUsuarioRequest {
  int32 id = 1;
  optional string nome = 2;
  optional int32 idade = 3;
}
\end{lstlisting}
\texttt{NOT\_FOUND}, \texttt{INVALID\_ARGUMENT} e \texttt{ALREADY\_EXISTS} correspondem aos três códigos do núcleo.
\end{frame}

\section{Comparação}

\begin{frame}{Semelhanças e diferenças}
  \centering
  \resizebox{\textwidth}{!}{%
  \begin{tabular}{@{}lllll@{}}
    \toprule
    & SOAP & REST & GraphQL & gRPC \\
    \midrule
    Contrato & WSDL & rotas e JSON & schema & \texttt{.proto} \\
    Mensagem & XML & JSON & JSON & binário \\
    Transporte & HTTP & HTTP & HTTP & HTTP/2 \\
    Estilo & operação & recurso & consulta & operação \\
    Cliente escolhe campos & não & não & sim & não \\
    Atualização parcial & campo opcional & campo omitido & argumento omitido & \texttt{optional} \\
    Inspeção manual & média & fácil & fácil & difícil \\
    \bottomrule
  \end{tabular}%
  }

  \vspace{0.6em}
  \small
  As quatro entregam o mesmo CRUD. A diferença prática apareceu no tamanho do contrato e em como a consulta de relacionamento é escrita.
\end{frame}

\begin{frame}{O que a implementação mostrou}
  \begin{itemize}
    \item Centralizar a regra em um módulo evitou quatro comportamentos diferentes para a cardinalidade 1..* da playlist.
    \item No REST, usuários, músicas e playlists caíram em coleções. As três consultas de ligação precisaram de rotas novas.
    \item No GraphQL, essas consultas são campos do schema, e o cliente monta o recorte. Foi o contrato mais curto para quem lê o pedido.
    \item SOAP e gRPC pediram uma operação por ação. O contrato fica explícito e cresce rápido: são cerca de vinte operações para três recursos.
    \item O gRPC foi o único que exigiu um passo de geração de código e campos \texttt{optional} para não tratar idade 0 como ``não enviado''.
    \item O SOAP foi o mais longo de escrever e o mais claro para uma ferramenta que lê WSDL.
  \end{itemize}
\end{frame}

\begin{frame}{Leitura crítica, ainda sem carga}
  Para este serviço, REST é a escolha mais direta: o domínio já é um conjunto de recursos. GraphQL compensa quando o cliente varia o recorte, como pedir a playlist com ou sem as músicas. SOAP se justifica quando o contrato formal e o ferramental pesam mais que o tamanho do XML. gRPC se justifica entre serviços, quando o contrato binário e o HTTP/2 importam mais que a facilidade de depurar.
  \vspace{0.6em}
  Essa leitura vem da experiência de implementar o mesmo CRUD. Ela ainda não diz qual tecnologia aguenta mais pedidos por segundo. Isso depende do teste de carga, que ficou de fora deste beta.
\end{frame}

\begin{frame}{Testes de carga}
  Esta versão não mede latência nem vazão.
  \vspace{0.5em}
  A versão final deve repetir as mesmas cinco consultas nas quatro tecnologias, com a mesma carga inicial e com concorrência crescente, e colocar os gráficos aqui. Os números não foram estimados de propósito: um gráfico sem medição atrapalharia a comparação.
\end{frame}

\begin{frame}{Conclusão do beta}
  \begin{itemize}
    \item O serviço cobre criação, consulta, alteração e remoção de usuários, músicas e playlists, com as ligações do diagrama.
    \item SOAP, REST, GraphQL e gRPC expõem essas operações em Python, sobre o mesmo núcleo.
    \item O próximo passo da entrega é o teste de carga e o preenchimento dos nomes da equipe.
  \end{itemize}
\end{frame}

\begin{frame}{Referências}
  \scriptsize
  \begin{enumerate}
    \item W3C. SOAP Version 1.2 Part 1: Messaging Framework. \url{https://www.w3.org/TR/soap12/}
    \item Fielding, R. T. Architectural Styles and the Design of Network-Based Software Architectures. UC Irvine, 2000.
    \item REST API Tutorial. \url{https://restfulapi.net/}
    \item GraphQL. \url{https://graphql.org/}
    \item gRPC. \url{https://grpc.io/}
    \item Stowe, M. XML, SOAP, JSON, REST, GraphQL?
    \item Brito, G.; Valente, M. T. REST vs GraphQL: A Controlled Experiment. ICSA, 2020.
  \end{enumerate}
\end{frame}

\end{document}
